% ABSTRACT (NBR 6028)
% --------------------------------------------------------------------------
\begin{resumo}[Abstract]
\begin{otherlanguage*}{english}
The democratization of access to Low Earth Orbit (LEO) via CubeSat-class nanosatellites imposes severe systems engineering challenges. During launch and deployment, the primary chassis is subjected to harsh dynamic environments, including transient shock inside the Poly-Picosatellite Orbital Deployer (P-POD) and random acoustic vibration governed by NASA GSFC-STD-7000A (GEVS). Conventional monolithic chassis exhibit high mass and limited dynamic isolation, transferring critical vibration levels directly to sensitive electronic payloads. This work presents a high-performance computational framework for the parametric modeling, physics-guided deep learning acceleration, and multi-objective optimization of a 1U CubeSat chassis featuring re-entrant auxetic metamaterial panels manufactured via Laser Powder Bed Fusion (L-PBF) in aerospace AlSi10Mg alloy. A 3D parametric CAD model incorporating strict Design for Additive Manufacturing (DfAM) gates was coupled with batch Ansys MAPDL finite element simulations (Block Lanczos prestressed modal extraction and random vibration PSD analysis), producing a high-fidelity dataset of 250 design points. A residual deep neural surrogate model (Physics-Guided ResNet) was developed with a composite loss function penalizing deviations from Gibson-Ashby scaling laws and enforcing stiffness monotonicity, achieving an average $R^2$ of $0.9874$ and an inference latency of \SI[output-decimal-marker={.}]{0.405}{\milli\second} (more than \num[group-separator={,}]{200000}$\times$ speedup over FEA). A deterministic state machine built with LangGraph as a strongly typed dataflow graph orchestrated early rejection of non-compliant geometries. Multi-objective evolutionary search via NSGA-II over \num[group-separator={,}]{10000} candidate evaluations mapped the Pareto frontier between structural mass and payload transmissibility. Using the TOPSIS multi-criteria decision method, the optimal flight design was selected ($\theta = \SI[output-decimal-marker={.}]{65.75}{\degree}$, $t = \SI[output-decimal-marker={.}]{0.613}{\milli\metre}$, $\nu_{\text{eff}} = -13.81$). High-order FEA ground truth verification showed close agreement with the neural predictions (\SI[output-decimal-marker={.}]{0.99}{\percent} error in $3\sigma$ peak stress and \SI[output-decimal-marker={.}]{0.008}{\percent} in transmissibility). Compared with the conventional solid aluminum chassis, the optimized auxetic design achieved a mass reduction of \SI[output-decimal-marker={.}]{62.49}{\percent} (\SI[output-decimal-marker={.}]{0.308}{\kilo\gram} $\to$ \SI[output-decimal-marker={.}]{0.116}{\kilo\gram}), passive elastodynamic vibration attenuation of \SI[output-decimal-marker={.}]{76.50}{\percent} (\SI[output-decimal-marker={.}]{12.6}{\decibel} reduction via mechanical wave filtering and acoustic impedance mismatch), a fundamental frequency of \SI[output-decimal-marker={.}]{579.62}{\hertz}, and full compliance with NASA GEVS stress margins and Steinberg random fatigue limits ($D = 0.046 \ll 0.25$), demonstrating flight qualification readiness.

\vspace{\onelineskip}
\noindent
\textbf{Keywords}: 1U CubeSat. Auxetic metamaterials. Design for additive manufacturing. Physics-guided neural networks. NSGA-II multi-objective optimization. NASA GEVS.
\end{otherlanguage*}
\end{resumo}

